
\UnnumberedChapter{KẾT LUẬN}

Trong khuôn khổ học phần Nhập môn Học máy,
nhóm đã thực hiện đề tài
\textbf{``Phân loại cảm xúc từ phản hồi của sinh viên''}
nhằm vận dụng các kiến thức về học máy có giám sát
và xử lý ngôn ngữ tự nhiên vào bài toán phân loại
văn bản.

Thông qua quá trình nghiên cứu và thực nghiệm,
nhóm đã xây dựng được một quy trình học máy
bao gồm các bước: khám phá dữ liệu (EDA),
tiền xử lý văn bản, biểu diễn đặc trưng,
huấn luyện mô hình, đánh giá kết quả và
phân tích các trường hợp dự đoán sai.

Bộ dữ liệu được sử dụng gồm 1.223 phản hồi
thuộc ba lớp cảm xúc: positive, neutral
và negative. Nhóm thực hiện phân chia dữ liệu
thành 978 mẫu huấn luyện và 245 mẫu kiểm thử
theo tỷ lệ 80/20, đồng thời áp dụng các
biện pháp hạn chế rò rỉ dữ liệu để bảo đảm
tính nhất quán trong quá trình thực nghiệm.

Về phương pháp, nhóm đã triển khai hai
cách biểu diễn văn bản là Bag-of-Words
và TF-IDF, kết hợp với hai thuật toán
Multinomial Naive Bayes và K-Nearest
Neighbors (K-NN). Các mô hình được
thử nghiệm với nhiều cấu hình n-gram
và siêu tham số khác nhau. Quá trình
lựa chọn mô hình được thực hiện bằng
Stratified 3-Fold Cross-Validation
với Macro F1-score làm tiêu chí đánh giá.

Kết quả thực nghiệm cho thấy mô hình
TF-IDF Unigram kết hợp K-NN với $k=5$
đạt Macro F1-score trung bình cao nhất
trên Cross-Validation, bằng 0,3647.
Tuy nhiên, khi đánh giá trên tập kiểm thử,
mô hình chỉ đạt Accuracy 0,3388 và
Macro F1-score 0,3369.

Trong khi đó, mô hình baseline sử dụng
Bag-of-Words Unigram kết hợp
Multinomial Naive Bayes đạt Macro
F1-score 0,3748 trên cùng tập kiểm thử.
Kết quả này cho thấy mô hình được lựa chọn
bằng Cross-Validation chưa đạt hiệu quả
tốt hơn baseline trong thực nghiệm hiện tại.

Bên cạnh việc so sánh các chỉ số đánh giá,
nhóm đã sử dụng ma trận nhầm lẫn và
phân tích các trường hợp dự đoán sai
để nhận diện những khó khăn của mô hình.
Qua đó, nhóm nhận thấy hiệu quả phân loại
có thể chịu ảnh hưởng bởi chất lượng nội dung
phản hồi, phương pháp biểu diễn đặc trưng
và phạm vi tham số được thử nghiệm.

Mặc dù kết quả dự đoán còn hạn chế và
chưa đủ cơ sở để khẳng định khả năng
ứng dụng trong môi trường thực tế,
đề tài đã đạt được mục tiêu xây dựng,
thử nghiệm và đánh giá một quy trình
học máy cơ bản theo yêu cầu học phần.

Thông qua quá trình thực hiện, nhóm có
cơ hội củng cố kiến thức về xử lý dữ liệu,
biểu diễn văn bản, lựa chọn thuật toán,
thiết kế thí nghiệm và diễn giải kết quả.
Đồng thời, việc tổ chức mã nguồn bằng
Jupyter Notebook và quản lý dự án trên
GitHub góp phần rèn luyện kỹ năng phối hợp
nhóm và phát triển phần mềm có khả năng
kiểm tra, tái tạo.

Trong tương lai, nhóm có thể tiếp tục
nghiên cứu cải thiện chất lượng dữ liệu,
mở rộng phạm vi thử nghiệm siêu tham số,
áp dụng các phương pháp biểu diễn văn bản
nâng cao và đánh giá mô hình trên những
tập dữ liệu độc lập khác.

Nhìn chung, đề tài không chỉ giúp nhóm
vận dụng kiến thức lý thuyết vào một
bài toán cụ thể mà còn làm rõ tầm quan
trọng của chất lượng dữ liệu, quy trình
thực nghiệm và phương pháp đánh giá
trong việc xây dựng các mô hình học máy.
